%==================================================================
\section{Syllabus map: what each class covered}

\begin{center}
\begin{longtable}{@{}p{1.4cm}p{13.4cm}@{}}
\toprule
\textbf{Class} & \textbf{Content} \\
\midrule
\endhead
\textbf{1} &
Motivation and vocabulary. Mahalanobis, jute production, the NSS (1950), IPNS, pilot surveys;
the Mahabharata anecdote. \textbf{The ABCD}: observation unit, target population, sampled
population, sampling unit, sampling frame. Census vs.\ sample survey. \textbf{The \emph{Literary
Digest} 1936 failure}. What sampling buys. \textbf{Total error $=$ sampling $+$ non-sampling}.
Probability vs.\ non-probability sampling; convenience, judgment, self-selection, snowball,
quota. \\[2pt]
\textbf{2} &
Finite-population framework: $U=\{1,\dots,N\}$, $Y_j$ \textbf{fixed}, parameters $\Ybar$,
$\sigma^2$. \textbf{Sampling design} $p:\mathcal S\to[0,1]$; \textbf{sampling scheme}. SRSWR and
SRSWOR. Implementation: lottery, random number tables (rejection and remainder methods),
\texttt{sample()}. SRSWR design $p(s)=1/N^n$, $\pi_j=1-(1-1/N)^n$. SRSWR estimation: $\ybar$
unbiased, $\Var(\ybar)=\sigma^2/n$, $\E(s^2)=\sigma^2$. SRSWOR design $p(s)=1/\binom Nn$,
$\pi_j=n/N$, $\pi_{jj'}=\frac{n(n-1)}{N(N-1)}$. \\[2pt]
\textbf{3} &
SRSWOR estimation: $\Var(\ybar)=\frac{N-n}{N-1}\frac{\sigma^2}{n}=(1-f)\frac{S^2}{n}$,
\textbf{finite population correction}; $\E(s^2)=S^2$. SRSWOR beats SRSWR for $n\ge2$; when SRSWR
is used (bootstrap). \textbf{Effective sample size} $n_{\text{eff}}$, $\ybar_{\text{eff}}$. Total,
variance, covariance. \textbf{Proportions}. Accuracy: SE, \textbf{CV/RSD}, Statistics Canada
guidelines. \\[2pt]
\textbf{4} &
\textbf{Confidence intervals}: why exact CIs are impossible; \textbf{H\'ajek (1960)}; the
approximate CI; $t_{n-1}$ for small $n$. \textbf{Choosing $n$}: margin of error, absolute $e$ and
relative $r$; $n_{WR}=z^2_{\alpha/2}\sigma^2/e^2$, $n_{WOR}\approx n_{WR}/(1+n_{WR}/N)$;
proportions and the $P(1-P)\le\frac14$ bound. \\[2pt]
\textbf{5} &
Worst case $P=0.5$ ($e=0.05\Rightarrow n=385$). Supplying $S^2$: worst-case, \textbf{pilot
survey}, previous studies, guess. Relative error $\Rightarrow$ the \textbf{CV}. \textbf{Testing
perspective}: $\alpha$, $\beta$, power, \textbf{effect size} $\Delta$;
$n_{WR}=(z_\alpha+z_\beta)^2\sigma^2/\Delta^2$; proportions. Worked: diabetes ($n=31$),
Framingham LDL ($n=809$). \\[2pt]
\textbf{6} &
\textbf{Stratified random sampling.} Why SRS can fail (representative only \emph{in
expectation}); auxiliary information. Setup: $H$ strata, $N_h$, $n_h$, $\sum N_h=N$,
$\sum n_h=n$. $\ybar_{st}=\frac1N\sum N_h\ybar_h$ is unbiased;
$\Var(\ybar_{st})=\frac1{N^2}\sum N_h^2(1-f_h)\frac{S_h^2}{n_h}$; variance estimation and CI.
\textbf{Allocation}: equal, proportional. \\[2pt]
\textbf{7} &
\textbf{Comparison of variances}: the exact identity for
$\Var_{SRS}(\ybar)-\Var_{st}(\ybar_{st})$ under proportional allocation; when stratification
helps (small WSS, large BSS). \textbf{Neyman allocation} $n_h\propto N_hS_h$, and that it
minimises $\Var(\ybar_{st})$ subject to $\sum n_h=n$. \textbf{Optimal allocation with a cost
constraint}: $n_h\propto N_hS_h/\sqrt{C_h}$. \\[2pt]
\textbf{8} &
Remarks on stratification (strata-level parameters, different frames and field procedures per
stratum, formation of strata). \textbf{Systematic sampling}: $N=nk$, random start in
$\{1,\dots,k\}$, sampling interval $k$; uses; behaviour under sorted vs.\ periodic frames.
Estimation: $\ybar$ unbiased, $\Var(\ybar)=\sigma_b^2$ (between-cluster).
\textbf{Intraclass correlation} $\rho_c$ and
$\Var(\ybar)=\frac{\sigma^2}{n}\{1+(n-1)\rho_c\}$. \\[2pt]
\textbf{9} &
Systematic sampling inference. \textbf{The variance cannot be estimated unbiasedly} from one
sample (a cluster-level sample of size 1). Pretending it is SRSWOR: $\hat V=(1-f)s^2/n$ is
\textbf{biased}, unbiased iff $\rho_c=-\frac1{N-1}$. Alternatives: $m$ independent systematic
samples (\textbf{IPNS}), multiple random starts, splitting one sample, Wolter (1984)
successive-difference estimator. \textbf{$N$ not a multiple of $n$}; \textbf{circular systematic
sampling}. \\
\bottomrule
\end{longtable}
\end{center}

\begin{note}{The results that carry the course}
\[
  \textbf{SRSWR:}\quad \E(\ybar)=\Ybar,\qquad \Var(\ybar)=\frac{\sigma^2}{n},\qquad \E(s^2)=\sigma^2 .
\]
\[
  \textbf{SRSWOR:}\quad \Var(\ybar)=\frac{N-n}{N-1}\cdot\frac{\sigma^2}{n}=(1-f)\frac{S^2}{n},
  \qquad \E(s^2)=S^2,\qquad S^2=\frac{N}{N-1}\sigma^2 .
\]
\[
  \textbf{Stratified:}\quad \ybar_{st}=\sum_h W_h\ybar_h,\qquad
  \Var(\ybar_{st})=\frac{1}{N^2}\sum_h N_h^2(1-f_h)\frac{S_h^2}{n_h},\qquad W_h=\frac{N_h}{N}.
\]
\[
  \textbf{Allocation:}\quad \text{prop. }n_h=nW_h;\quad \text{Neyman }n_h=\frac{nN_hS_h}{\sum N_{h'}S_{h'}};\quad
  \text{cost }n_h\propto\frac{N_hS_h}{\sqrt{C_h}} .
\]
\[
  \textbf{Systematic:}\quad \Var(\ybar)=\sigma_b^2=\frac{\sigma^2}{n}\{1+(n-1)\rho_c\},
  \qquad -\frac{1}{n-1}\le\rho_c\le1 .
\]
\[
  \textbf{Sample size:}\quad n_{WR}=\frac{z^2_{\alpha/2}\sigma^2}{e^2}\ \text{or}\
  \frac{(z_\alpha+z_\beta)^2\sigma^2}{\Delta^2},\qquad
  n_{WOR}\approx\frac{n_{WR}}{1+n_{WR}/N}.
\]
\end{note}

\begin{trap}{the three ``opposite'' facts that get mixed up}
\begin{itemize}
  \item \textbf{Stratified} sampling wants strata \emph{internally homogeneous} (small $S_h$) and
        strata means far apart. \textbf{Systematic} sampling wants clusters \emph{internally
        heterogeneous} (small $\rho_c$) and clusters resembling one another. \textbf{Exactly
        opposite.}
  \item $s^2$ is unbiased for $\sigma^2$ under SRSWR but for $S^2$ under SRSWOR. Same statistic,
        different target.
  \item Under systematic sampling $\ybar$ \emph{is} unbiased, but its \emph{variance} cannot be
        estimated unbiasedly at all --- a distinction worth stating explicitly.
\end{itemize}
\end{trap}

%==================================================================
\section{Concepts and vocabulary}

\Q{}{5}
Define, and illustrate each with the running example \emph{``estimate the average monthly
out-of-pocket medical expenditure of all adults in West Bengal''}: (i) observation unit,
(ii) target population, (iii) sampled population, (iv) sampling unit, (v) sampling frame.
Explain what goes wrong when the sampled population differs from the target population.

\Q{}{5}
Distinguish a \emph{census} from a \emph{sample survey}. Give three distinct reasons why a sample
survey can be \textbf{more accurate} than a complete enumeration.

\Q{}{8}
In 1936 the \emph{Literary Digest} mailed 10 million questionnaires, received 2.3 million
replies, and predicted Landon would beat Roosevelt. Identify the \textbf{two distinct} sources of
error, explain the mechanism of each, and explain precisely why 2.3 million responses did not
rescue the prediction. What does this say about the relation between sample size and bias?

\Q{}{5}
Write down the decomposition \emph{total error $=$ sampling error $+$ non-sampling error}. Define
each, give two examples of non-sampling error, and explain why increasing $n$ can leave the total
error unchanged or increase it.

\Q{}{8}
Define probability sampling. For each of \emph{convenience}, \emph{judgment}, \emph{self-selection},
\emph{snowball} and \emph{quota} sampling give a one-line description, one situation where it is
the pragmatic choice, and the specific bias it is most vulnerable to.

\Q{}{8}
Carefully distinguish a \textbf{sampling design} from a \textbf{sampling scheme}. Give an example
of two genuinely different schemes inducing the \emph{same} design, and explain why the design is
what inference depends on.

%==================================================================
\section{Designs, selection and inclusion probabilities}

\Q{}{8}
Let $U=\{1,2,3\}$ and $n=2$.
\pt{a} Write down $\mathcal S$ and the SRSWOR design, and verify $\sum_sp(s)=1$.
\pt{b} Compute $\pi_j$ and $\pi_{jj'}$ directly and check against the general formulas.
\pt{c} Verify $\sum_j\pi_j=n$ and explain why this holds for \emph{any} fixed-size design.

\Q{}{10}
\pt{a} Show SRSWR induces the uniform design $p(\bm s)=1/N^n$ on \emph{ordered} samples.
\pt{b} Derive $\pi_j=1-\left(1-\frac1N\right)^n$.
\pt{c} Show $\pi_j<n/N$ for $n\ge2$ and interpret the gap.
\pt{d} What is $\lim_{n\to\infty}\pi_j$, and what does this permit that SRSWOR does not?

\Q{}{10}
\pt{a} Show SRSWOR induces the uniform design $p(\bm s)=1/\binom Nn$ on \emph{unordered} samples.
\pt{b} Derive $\pi_j=n/N$ and $\pi_{jj'}=\frac{n(n-1)}{N(N-1)}$.
\pt{c} Verify $\sum_{j\ne j'}\pi_{jj'}=n(n-1)$.

\Q{}{8}
Under SRSWOR show that the \emph{marginal} probability that unit $j$ is selected at the $k$-th
draw equals $1/N$ for every $k$. \src{This is the ``why?'' on the Class-2 slide.}

\Q{}{5}
With $N=17$ and $n=4$, describe the \textbf{rejection} and \textbf{remainder} methods for drawing
an SRSWR from a table of random digits, carry out one draw of each from a stated row, and compare
their rejection rates.

%==================================================================
\section{Estimation under SRSWR and SRSWOR}

\Q{}{10}
Under SRSWR: \pt{a} show $\ybar$ is unbiased for $\Ybar$; \pt{b} show $\Var(\ybar)=\sigma^2/n$;
\pt{c} state precisely where the ``with replacement'' assumption is used in each part, and why
$\Var(\ybar)$ does not depend on $N$ except through $\sigma^2$.

\Q{}{8}
Show $\E(s^2)=\sigma^2$ under SRSWR, and deduce that $\wh\Var(\ybar)=s^2/n$ is unbiased.

\Q{}{8}
For any known $h(\cdot)$, show $\frac1n\sum_ih(y_i)$ is unbiased for $\frac1N\sum_jh(Y_j)$ under
SRSWR. Explain why this needs \textbf{no distributional assumption}, and what exactly makes it
true.

\Q{}{12}
Under SRSWOR, with $y_1,\dots,y_n$ the values at the $n$ draws:
\pt{a} Show $\E(y_i)=\Ybar$, $\Var(y_i)=\sigma^2$.
\pt{b} Show $\Cov(y_i,y_{i'})=-\dfrac{\sigma^2}{N-1}$ for $i\ne i'$.
\pt{c} Derive $\Var(\ybar)=\dfrac{N-n}{N-1}\cdot\dfrac{\sigma^2}{n}=(1-f)\dfrac{S^2}{n}$.
\pt{d} Interpret the sign of the covariance and why it makes SRSWOR more efficient.

\Q{}{10}
Show $\E(s^2)=S^2$ under SRSWOR and deduce that $\wh\Var(\ybar)=(1-f)s^2/n$ is unbiased.
\src{Note the target is $S^2$, not $\sigma^2$.}

\Q{}{8}
\pt{a} Prove $\Var_{WOR}(\ybar)<\Var_{WR}(\ybar)$ for $n\ge2$, with equality iff $n=1$.
\pt{b} Interpret the fpc $(1-f)$ and describe the regimes $f\to0$ and $f\to1$.
\pt{c} Give three reasons SRSWR is nonetheless used in practice.

\Q{}{8}
For the population total $T=\sum_jY_j$ and $\wh T=N\ybar$:
\pt{a} Show $\wh T$ is unbiased and find $\Var(\wh T)$ under SRSWOR.
\pt{b} Give an unbiased estimator of $\Var(\wh T)$.
\pt{c} Show $\CV(\wh T)=\CV(\ybar)$ and say why that is useful in a survey report.

%==================================================================
\section{Effective sample size and proportions}

\Q{}{10}
Under SRSWR let $n_{\text{eff}}$ be the number of \emph{distinct} units sampled.
\pt{a} Explain why it is random.
\pt{b} Show $\E(n_{\text{eff}})=N\left[1-\left(1-\frac1N\right)^n\right]$.
\pt{c} Evaluate for $N=100$, $n=20$ and comment.
\pt{d} Show $\E(n_{\text{eff}})\le n$ and say when the two are close.

\Q{}{8}
Let $\ybar_{\text{eff}}$ be the mean of $Y$ over the distinct units of an SRSWR sample. Show
$\E(\ybar_{\text{eff}})=\Ybar$. \src{Hint: condition on the set of distinct units and use
symmetry.} Compare $\ybar_{\text{eff}}$ with the SRSWOR mean of size $n$.

\Q{}{10}
Let $Y_j=\I(\text{unit }j\text{ has the characteristic})$ and $P$ the population proportion.
\pt{a} Show $P=\Ybar$.
\pt{b} Show $\sigma^2=P(1-P)$ and $S^2=\frac{N}{N-1}P(1-P)$.
\pt{c} Write $\Var(p)$ under SRSWR and SRSWOR.
\pt{d} Show $\wh\Var(p)=(1-f)\dfrac{p(1-p)}{n-1}$ is unbiased under SRSWOR.

\Q{}{8}
\pt{a} Show $np\sim\text{Binomial}(n,P)$ under SRSWR, and state the distribution under SRSWOR.
\pt{b} Using these, identify the exact factor by which the SRSWOR variance is smaller.

%==================================================================
\section{Accuracy, confidence intervals, sample size}

\Q{}{5}
Define the standard error and the coefficient of variation of an estimator, give
$\wh\SE(\ybar)$ and $\wh\CV(\ybar)$ under SRSWOR, and state the Statistics Canada guidelines.

\Q{}{10}
An SRSWOR with $N=5{,}000$, $n=100$ gives $\ybar=\text{Rs.}\,2{,}400$, $s=\text{Rs.}\,1{,}800$.
\pt{a} Compute $\wh\Var(\ybar)$, $\wh\SE(\ybar)$, $\wh\CV(\ybar)$.
\pt{b} Classify by the Statistics Canada guidelines.
\pt{c} Give an approximate 95\% CI for $\Ybar$.
\pt{d} Estimate the total and give its estimated CV.
\pt{e} By what factor must $n$ increase to halve the CV?

\Q{}{10}
\pt{a} Explain why an \emph{exact} 95\% CI for $\Ybar$ cannot be computed in finite-population
sampling, even though the sampling distribution of $\ybar$ is fully determined by the design.
\pt{b} State H\'ajek's (1960) result and its conditions.
\pt{c} Write the resulting approximate CI under SRSWOR and say when $t_{n-1}$ should replace
$1.96$.

\Q{}{8}
A survey reports a 95\% CI of $[4735,7265]$ for mean monthly rent, from $n=15$.
\pt{a} State precisely what the coverage probability refers to.
\pt{b} A student says ``there is a 95\% probability that $\Ybar$ lies in this interval''. Explain
why that is wrong here, and give the correct statement.
\pt{c} Comment on the interval's usefulness and on what should have been done instead.

\Q{}{12}
\pt{a} From $\Prob(|\ybar-\Ybar|\le e)=1-\alpha$ and the normal approximation, derive
$n_{WR}=\frac{z^2_{\alpha/2}\sigma^2}{e^2}$.
\pt{b} Derive $n_{WOR}=\frac{z^2_{\alpha/2}S^2}{e^2+z^2_{\alpha/2}S^2/N}$.
\pt{c} Hence show $n_{WOR}=\frac{n_{WR}}{1+n_{WR}/N}$, and deduce $n_{WOR}<N$ always.
\pt{d} Comment on the claim that one should sample ``a fixed percentage of the population''.

\Q{}{10}
\pt{a} Derive $n_{WR}=\frac{z^2_{\alpha/2}P(1-P)}{e^2}$ for a proportion.
\pt{b} Show $P(1-P)\le\frac14$ and obtain the worst-case $n$.
\pt{c} At 95\% confidence compute the worst-case $n$ for $e=0.10,\,0.05,\,0.03$ and comment on
the scaling.
\pt{d} If in truth $P=0.2$, quantify how wasteful the worst-case choice was.

\Q{}{8}
\pt{a} Show that with a \emph{relative} margin of error $r$, the required $n$ depends on the
population \textbf{coefficient of variation}. Derive the formula.
\pt{b} Explain why the CV is easier to guess in advance than $\sigma^2$.

\Q{}{12}
For $H_0:\Ybar=\mu_0$ against $H_1:\Ybar=\mu_1>\mu_0$, with $\Delta=\mu_1-\mu_0$:
\pt{a} Write the test statistic and level-$\alpha$ rejection rule, stating the approximation used.
\pt{b} Derive $n_{WR}=\frac{(z_\alpha+z_\beta)^2\sigma^2}{\Delta^2}$ from the power requirement.
\pt{c} State the two-sided and SRSWOR versions.
\pt{d} Explain the role of the \emph{minimum detectable effect} and why $\Delta$ is a
subject-matter choice.

\Q{}{10}
Among people free of diabetes, mean fasting glucose is $95.0$ mg/dL with sd $9.8$; a mean of
$100$ in coffee drinkers would be clinically important. Under SRSWR, how many patients give a
two-sided 5\% test power 80\%? Show all working and state assumptions.
\src{$z_{0.025}=1.96$, $z_{0.20}=0.8416$.}

\Q{}{10}
26\% of people free of heart disease have elevated LDL. How many patients are needed for a
two-sided 5\% test to have power 90\% to detect a difference of $0.05$ in the proportion?
\src{$z_{0.025}=1.96$, $z_{0.10}=1.2816$.}

\Q{}{8}
Your required $n$ is three times the budget. Describe three distinct ways to reduce it, state
exactly what each costs you, and say which you would choose and why.

%==================================================================
\section{Stratified random sampling}

\Q{}{8}
A campus health centre wants the mean height of $1{,}000$ undergraduates, half male and half
female. An SRSWOR of size $10$ happens to return $9$ men and $1$ woman.
\pt{a} Is the estimator $\ybar$ biased? Answer precisely, and explain the difference between
being representative \emph{in expectation} and being representative \emph{in this sample}.
\pt{b} Explain how stratification fixes the problem, and what extra information it requires.
\pt{c} Is a stratified sample still a probability sample? Justify.

\Q{}{12}
Set up stratified random sampling with $H$ strata, $N_h$, $n_h$, and SRSWOR drawn
\textbf{independently} within each stratum.
\pt{a} Show that $\ybar=\frac1n\sum\text{(all sampled values)}$ is in general \emph{biased} for
$\Ybar$, and state exactly when it is unbiased.
\pt{b} Show that $\ybar_{st}=\frac1N\sum_hN_h\ybar_h$ is unbiased.
\pt{c} Derive $\Var(\ybar_{st})=\frac1{N^2}\sum_hN_h^2(1-f_h)\frac{S_h^2}{n_h}$, stating clearly
where the independence across strata is used.
\pt{d} Write down an unbiased estimator of this variance and the approximate 95\% CI for $\Ybar$.

\Q{}{10}
\pt{a} Define equal, proportional and Neyman allocation, and state the one-line rationale for
each.
\pt{b} Show that under \textbf{proportional} allocation $f_h=f$ for every $h$, and hence
$\Var_{prop}(\ybar_{st})=\frac{1-f}{n}\sum_hW_hS_h^2$.
\pt{c} Show that under \textbf{Neyman} allocation
$\Var_{Ney}(\ybar_{st})=\frac{1}{n}\left(\sum_hW_hS_h\right)^2-\frac1N\sum_hW_hS_h^2$.
\pt{d} Deduce $\Var_{prop}\ge\Var_{Ney}$, and state exactly when equality holds.

\Q{}{10}
Prove that Neyman allocation minimises $\Var(\ybar_{st})$ subject to $\sum_hn_h=n$.
\src{Use a Lagrange multiplier, or the Cauchy--Schwarz inequality.}

\Q{}{10}
Suppose the cost of the survey is $A+\sum_hn_hC_h$, where $A$ is overhead and $C_h$ the cost per
unit in stratum $h$.
\pt{a} Formulate the problem of minimising $\Var(\ybar_{st})$ subject to a total budget $C$.
\pt{b} Show the optimum is $n_h\propto\dfrac{N_hS_h}{\sqrt{C_h}}$, and write $n_h$ explicitly.
\pt{c} Show this reduces to Neyman allocation when $C_h$ is constant, and explain intuitively why
expensive strata get smaller samples.

\Q{}{12}
Compare an SRSWOR of size $n$ with a stratified WOR sample of the same size $n$ under
proportional allocation. The exact identity is
\[
  \Var_{SRS}(\ybar)-\Var_{st}(\ybar_{st})
  =\frac{N-n}{n(N-1)}\left[\sum_{h=1}^HW_h(\Ybar_h-\Ybar)^2-\frac1N\sum_{h=1}^H(1-W_h)S_h^2\right].
\]
\pt{a} Identify the two bracketed terms as between-stratum and within-stratum contributions.
\pt{b} State the condition under which stratification helps, in words.
\pt{c} What happens in the special case $\Ybar_1=\cdots=\Ybar_H$? What has gone wrong, and what
does it tell you about how to form strata?
\pt{d} Why is stratification nevertheless ``almost always beneficial'' in realistic surveys?

\Q{}{12}
A district has $N=1{,}000$ households in three strata: urban ($N_1=200$), semi-urban
($N_2=300$), rural ($N_3=500$). From a pilot, $S_1=60$, $S_2=40$, $S_3=20$ and
$\Ybar_1=300$, $\Ybar_2=200$, $\Ybar_3=100$ (rupees per month). Take $n=100$.
\pt{a} Compute the proportional and the Neyman allocations.
\pt{b} Compute $\Var_{prop}(\ybar_{st})$ and $\Var_{Ney}(\ybar_{st})$.
\pt{c} Compute $\Ybar$ and (using $S^2\approx\sum W_hS_h^2+\sum W_h(\Ybar_h-\Ybar)^2$) the
variance under SRSWOR of the same size. Compare all three and comment.
\pt{d} Neyman allocation gives non-integer $n_h$. How would you round, and what problem does
rounding create?

\Q{}{8}
Give four practical reasons, beyond variance reduction, for using stratified sampling. For each,
give a concrete example. \src{Class 8, ``Stratification: some remarks''.}

%==================================================================
\section{Systematic sampling}

\Q{}{10}
Let $N=nk$ with $k$ an integer.
\pt{a} State the linear systematic sampling scheme precisely.
\pt{b} Write down $\mathcal S$, the sampling design, and show $\pi_j=n/N$ for every $j$.
\pt{c} Compute $\pi_{jj'}$ and show it is $0$ for some pairs. Explain why this single fact is
responsible for the variance-estimation problem of Class 9.

\Q{}{12}
With $N=nk$, write $Y_{ij}$ for the $j$-th unit of cluster $i$, $i=1,\dots,k$.
\pt{a} Show that a systematic sample is exactly one randomly chosen cluster, and hence a special
case of cluster sampling.
\pt{b} Show $\E(\ybar)=\Ybar$.
\pt{c} Show $\Var(\ybar)=\frac1k\sum_{i=1}^k(\Ybar_i-\Ybar)^2=\sigma_b^2$, the
\textbf{between-cluster} variance.
\pt{d} Deduce the design rule: for systematic sampling to be efficient, should clusters be
internally homogeneous or heterogeneous? Contrast with stratified sampling.

\Q{}{10}
Let $\rho_c$ be the intraclass correlation coefficient.
\pt{a} Using $\Var(\ybar)=\frac{\sigma^2}{n}\{1+(n-1)\rho_c\}$, show that systematic sampling
beats SRSWR iff $\rho_c\le0$.
\pt{b} Show that it beats SRSWOR iff $\rho_c\le-\dfrac1{N-1}$.
\pt{c} Show $-\dfrac1{n-1}\le\rho_c\le1$ and interpret both extremes.

\Q{}{12}
Take $N=12$, $n=3$, $k=4$.
\pt{a} Population A has $Y=(1,2,\dots,12)$ in that order. List the four possible systematic
samples, compute $\Var(\ybar)$ directly, and compare with the SRSWR and SRSWOR variances.
\pt{b} Population B has $Y=(10,20,30,40)$ repeated three times. Repeat the computation.
\pt{c} Compute $\rho_c$ in each case from
$\Var(\ybar)=\frac{\sigma^2}{n}\{1+(n-1)\rho_c\}$, and explain the two results in terms of the
ordering of the sampling frame.

\Q{}{10}
\pt{a} Explain why $\Var(\ybar)$ cannot be estimated unbiasedly from a single systematic sample.
\pt{b} Suppose we pretend the sample is SRSWOR and use $\hat V=(1-f)\frac{s^2}{n}$. Given that
$\E(\hat V)=(1-f)\frac{S^2}{n}\cdot\frac{N-1}{N}(1-\rho_c)$, determine the direction of the bias
as a function of $\rho_c$, and state when $\hat V$ is unbiased.
\pt{c} When is this pretence a \emph{conservative} (safe) thing to do?

\Q{}{10}
Describe three practical ways of obtaining a usable variance estimate under systematic sampling,
and state the cost or assumption each one carries:
(i) $m$ independent systematic samples (IPNS);
(ii) $m$ random starts within the first $mk$ units;
(iii) the successive-difference estimator of Wolter (1984).
Write down the estimator in each case.

\Q{}{8}
The ISI canteen has $N=14$ staff and you want $n=4$.
\pt{a} Show that linear systematic sampling breaks down, and describe the two ad hoc fixes with
$k=3$ and $k=4$, listing the resulting samples and sizes.
\pt{b} Describe \textbf{circular systematic sampling} and carry it out with $k=4$ and random
start $i=13$.
\pt{c} Show that circular systematic sampling gives $\pi_j=n/N$ for every $j$, and explain why
this is an improvement on the ad hoc fixes.

\Q{}{8}
State four situations in which systematic sampling is preferred to SRS in practice, and two
situations in which it can go badly wrong. Give a concrete example of each failure.

%==================================================================
\section{Applied and mixed problems}

\Q{}{12}
A district has $N=1{,}200$ registered small businesses. An SRSWOR of $n=150$ finds that $42$
rely on local moneylenders rather than bank loans.
\pt{a} Estimate $P$ and give the estimated standard error.
\pt{b} Give an approximate 95\% CI for $P$.
\pt{c} Estimate the total number of such businesses, with its standard error.
\pt{d} Compute the estimated CV and classify it.
\pt{e} A follow-up wants margin of error $0.03$ at 95\%. Find the required $n$ (i) using the
worst-case bound and (ii) using this survey as a pilot.

\Q{}{10}
Critically evaluate the following design, identifying every flaw and classifying each as a
sampling or a non-sampling error.
\begin{quote}\small
\emph{To estimate the average monthly household expenditure on private tuition in a city, a
researcher stands outside three well-known coaching centres on a weekday evening and interviews
the first 200 parents who agree to talk. Of the 200 approached, 120 respond. The questionnaire
asks ``How much do you spend on your child's tuition?'' The researcher reports the sample mean as
the city average, with standard error $s/\sqrt{120}$.}
\end{quote}

\Q{}{10}
A population of $N=6$ units has $Y=(2,4,4,6,8,12)$.
\pt{a} Compute $\Ybar$, $\sigma^2$, $S^2$ and verify $S^2=\frac{N}{N-1}\sigma^2$.
\pt{b} For $n=2$ under SRSWOR, list all $\binom62$ samples, verify $\E(\ybar)=\Ybar$ directly and
check $\Var(\ybar)$ against $(1-f)S^2/n$.
\pt{c} Compute $\Var_{WR}(\ybar)$ and confirm the inequality.

\Q{}{12}
You must design a survey to estimate mean household electricity consumption in a city of
$N=100{,}000$ households. You know the ward of every household, and previous studies suggest
consumption varies far more \emph{between} wards than \emph{within} them. The budget allows
$n=2{,}000$.
\pt{a} Which of SRSWOR, stratified and systematic sampling would you choose, and why? Address all
three.
\pt{b} If you stratify by ward, which allocation would you use, and what would you need to know?
\pt{c} Suppose the household list is sorted by ward. What does a systematic sample of every
$50$th household resemble, and is that good or bad here?
\pt{d} State one non-sampling error this survey is particularly exposed to, and how you would
mitigate it.
